\section{Dynamics and time heterogeneity}\label{app:dynamics}

\begin{figure}[htbp]
\centering
\includegraphics[width=0.90\textwidth]{appendix_event_study.pdf}
\caption{Categorical beta Q5--Q1 event-time profile. The plotted object is the fitted tail contrast, not mapping coverage or continuous residual information. December 2022 is the omitted transition month; October 2025 is absent because the survey was not collected.}
\label{fig:app-event}
\end{figure}

The preperiod max-statistic is 1.502 with common-multiplier $p=.929$, so the implemented event study does not detect a differential pretrend. This does not establish parallel trends or attribute a later break to AI. Forty of 42 observed post coefficients are negative, but the path is neither immediate nor monotone.

\begin{table}[htbp]
\centering
\caption{Joint post-period coefficients}\label{tab:apptime}
\small
\begin{tabular}{lrrr}
\toprule
Period or contrast & Coefficient & Paired SE & 95\% interval\\
\midrule
2023&-.1356&.0452&[-.2245,-.0467]\\
2024&-.0783&.0548&[-.1821,.0254]\\
2025--26&-.1629&.0507&[-.2623,-.0635]\\
2024 minus 2023& .0573&.0447&[-.0302,.1447]\\
2025--26 minus 2023&-.0273&.0440&[-.1154,.0608]\\
2025--26 minus 2024&-.0846&.0374&[-.1577,-.0115]\\
\bottomrule
\end{tabular}
\end{table}

The post-2025 extension is different from 2024 but not from 2023. Interest rates, technology-sector adjustment, return-to-office changes, post-pandemic normalization, and CPS composition remain concurrent explanations. The failed 5,148-parameter occupation-age seasonality model is recorded rather than replaced.
\FloatBarrier
